%% kapa: exam=1 round=20 year=2009 period=2 form=A subject=appraisal_law subject_ko=부동산관계법규 seq=28 number=68
%% source: Q-Net 감정평가사 기출문제 (cst00309) artlSeq=5024178 file=extracted/1008/5024178/2009년 제20회 감정평가사 1차시험 기출문제/2교시.hwp
%% conversion: kordoc 4.24.51 -> tools/gen1.py (auto); review=proofread-sonnet; flags=none
\begin{question}{68}
국유재산법령상 무주부동산의 처리에 관한 설명으로 옳지 \uline{않은} 것은?
\begin{choices}
  \item 국유재산사무의 총괄청 또는 관리청은 무주의 부동산을 국유재산으로 취득한다.
  \item 무주의 부동산을 국유재산으로 취득함에 있어서는 대통령령이 정하는 바에 따라 6월 이상의 기간을 정하여 그 기간 내에 정당한 권리자 기타 이해관계인이 이의를 제기할 수 있다는 뜻을 공고하여야 한다.
  \item 무주부동산을 국유재산으로 취득하여 소유자 등록을 하는 경우 그 권리자의 명의는 국(國)으로 한다.
  \item 취득한 국유재산은 당해 국유재산이 공익사업을 위한 토지 등의 취득 및 보상에 관한 법률에 의한 공익사업에 필요하게 된 경우 그 취득일부터 10년간은 이를 매각할 수 없다.
  \item 부동산의 소유자가 따로 있음을 알 수 있는 부동산에 대하여 국가가 무주부동산 공고절차를 거쳐 국유재산으로 등기를 마치고 점유를 개시하였다면, 그 점유의 개시에 있어 자기의 소유라고 믿은 데 과실이 있다는 것이 판례의 입장이다.
\end{choices}
\end{question}
\answer{4}
